\originalchapter{平衡附近的振动}\label{ch:osc}
\originalsection{回复力与小振动}
轻拉吉他弦后放手, 弦反复经过静止时的位置. 门上的弹簧、钟摆和电路中的电荷也会反复变化. 这些现象共有的结构是: 系统偏离稳定平衡时受到使它返回的作用, 惯性又使它越过平衡. 在写方程以前, 必须确定哪个量在变化, 哪些力实际作用于它, 以及近似的适用范围.

以质量 $m$ 的物块沿水平无摩擦导轨运动为例. 坐标 $x$ 是相对平衡位置的位移, 不是弹簧总长度. 重力与支持力在竖直方向相抵; 若弹簧在考察范围内满足胡克关系, 水平合力为 $-kx$, 其中 $k>0$ 是弹簧劲度系数, 单位为 $\mathrm{N/m}$. 牛顿定律成为 $m\ddot x=-kx$. 定义固有角频率 $\omega_0=\sqrt{k/m}$, 则
\begin{equation}\label{eq:sho}
\ddot x+\omega_0^2x=0,\qquad
x(t)=x_0\cos\omega_0t+\frac{v_0}{\omega_0}\sin\omega_0t.
\end{equation}
这里 $x_0=x(0)$, $v_0=\dot x(0)$ 是准备系统时给定的初始条件. 直接求两次导数可验证此式满足方程, 代入 $t=0$ 可核对两初值. 为说明没有遗漏另一种运动, 若两个解有相同初值, 它们之差 $z$ 满足同一个齐次方程及零初值; $\dot z^2+\omega_0^2z^2$ 的导数为零, 初值为零, 故 $z=0$.

角频率以每秒弧度计, 普通频率 $f=\omega_0/(2\pi)$ 以赫兹计, 周期 $T_0=2\pi/\omega_0$. 不能把 $\omega_0$ 的数值直接当作赫兹数. 系统的总能量为 $E=m\dot x^2/2+kx^2/2$, 并且 $\dot E=\dot x(m\ddot x+kx)=0$. 动能与弹性势能不断交换, 总能量不变.

理想弹簧不是所有系统的本来面貌. 设一维保守系统的势能 $V(x)$ 在平衡点 $x_*$ 附近光滑, 写 $q=x-x_*$. 平衡要求 $V'(x_*)=0$. 泰勒展开给出
\[
V(x_*+q)=V(x_*)+\tfrac12 V''(x_*)q^2+\tfrac16 V'''(x_*)q^3+\cdots,
\quad F=-V''(x_*)q-\tfrac12 V'''(x_*)q^2+\cdots.
\]
当 $V''(x_*)>0$ 且被忽略项相对线性项很小时, $\omega_0^2=V''(x_*)/m$. 若三次导数非零, 可用 $|qV'''(x_*)|/(2V''(x_*))\ll1$ 检查首个修正; 还须检查更高次项. 若 $V''(x_*)=0$, 即使平衡稳定也不能套用非零频率的简谐模型. 例如 $V=aq^4$ 且 $a>0$ 的回复力为 $-4aq^3$.

\begin{example}
长 $\ell$ 的均匀细杆绕上端在竖直平面内摆动, 质量为 $m$. 求小振动的角频率, 并解释为何与点质量摆不同.
\end{example}
\begin{solution}
取 $\theta=0$ 为杆竖直向下, 正角方向与正角加速度一致. 重力矩为 $-mg(\ell/2)\sin\theta$, 转动惯量为 $I=m\ell^2/3$. 因而 $I\ddot\theta=-mg\ell\sin\theta/2$. 在 $|\theta|\ll1$ 弧度时, $\sin\theta=\theta+O(\theta^3)$, 得 $\omega_0=\sqrt{3g/(2\ell)}$. 点质量摆的全部质量在距离 $\ell$ 处, $I=m\ell^2$, 重力矩的力臂也为 $\ell$; 所以它的角频率为 $\sqrt{g/\ell}$. 回复作用与惯性必须同时比较.
\end{solution}
\originalsection{幅度、相位与复表示}
式 \eqref{eq:sho} 可写成 $x=A\cos(\omega_0t+\phi)$, 其中 $A\ge0$ 是幅度, $\phi$ 是初相位. 对 $x=a\cos\omega_0t+b\sin\omega_0t$, 有 $A=\sqrt{a^2+b^2}$, $a=A\cos\phi$, $b=-A\sin\phi$. 只使用 $\arctan(-b/a)$ 会丢失象限; 应同时依据正弦、余弦确定相位. $A=0$ 时相位没有物理意义.

复数的作用是统一处理幅度与相位. 采用 $x(t)=\Rea[\widetilde x\e^{\ii\omega t}]$, $\widetilde x=A\e^{\ii\phi}$; 微分对应乘 $\ii\omega$. 对实系数线性方程, 复解的实部仍是解. 这不意味着任何非线性物理量都能取实部: $x^2$ 通常不等于 $\Rea[(\widetilde x\e^{\ii\omega t})^2]$. 能量必须由实位移和实速度计算, 或使用正确的时间平均公式.

叠加原理来自方程的线性: 若 $x_1,x_2$ 是齐次方程的解, 常数 $c_1x_1+c_2x_2$ 也是解. 时间平移不变性则来自系数不依赖时间: $x(t+t_*)$ 仍满足同一方程. 它与叠加是两个不同性质; 例如时间变化的 $k(t)$ 保持线性, 却一般破坏时间平移不变性.

\originalsection{电路中的同一结构}
电感 $L$ 与电容 $C$ 构成无电阻的串联闭合回路. 选定电容极板上的电荷 $Q$ 和电流方向使 $I=\dot Q$, 基尔霍夫电压定律给出 $L\ddot Q+Q/C=0$. 此处电感对应惯性, 电容倒数对应回复系数; $\omega_0=1/\sqrt{LC}$. 磁场能 $LI^2/2$ 与电场能 $Q^2/(2C)$ 交换. 若 $Q(0)=0$, $I(0)=I_0$, 则 $Q=I_0\sin\omega_0t/\omega_0$. 回路箭头与电荷极性不同的图中, $I$ 可能是 $-\dot Q$, 必须先核对符号再代初值.
